\documentclass[10pt,a4paper]{article}
\usepackage[utf8]{inputenc}
\usepackage[T1]{fontenc}
\usepackage[ngerman]{babel}
\usepackage{geometry}
\usepackage{xcolor}
\usepackage{tabularx}
\usepackage{array}
\usepackage{booktabs}
\usepackage{enumitem}
\usepackage{helvet}
\usepackage{tikz}
\usepackage{colortbl}
\usepackage{amssymb}
\usepackage{fancyhdr}
\usepackage{paracol}
\usepackage{titlesec}
\usepackage{footmisc}

\renewcommand{\familydefault}{\sfdefault}
\geometry{margin=1.5cm, top=2cm, bottom=2cm}
\setlength{\parindent}{0pt}
\setlength{\parskip}{4pt}
\setlength{\headheight}{14pt}

% Colors
\definecolor{spanisch}{HTML}{c41e3a}
\definecolor{lightred}{HTML}{fdf2f4}
\definecolor{darkgray}{HTML}{333333}
\definecolor{lightgray}{HTML}{f5f5f5}
\definecolor{english}{HTML}{1a5276}
\definecolor{spaet}{HTML}{8e44ad}
\definecolor{lightspaet}{HTML}{f5eef8}

% Headers
\pagestyle{fancy}
\fancyhf{}
\fancyhead[L]{\small\textcolor{spaet}{\textbf{Sequenzplan Spätbeginner Spanisch M-V}}}
\fancyhead[R]{\small\textcolor{english}{\textbf{Late Starter Spanish Sequence Plan M-V}}}
\fancyfoot[C]{\thepage}

% Section formatting
\titleformat{\section}{\large\bfseries\color{spaet}}{}{0pt}{}
\titleformat{\subsection}{\normalsize\bfseries\color{darkgray}}{}{0pt}{}
\titlespacing{\section}{0pt}{10pt}{5pt}
\titlespacing{\subsection}{0pt}{6pt}{3pt}

% Section bilingual
\newcommand{\secbil}[2]{%
\vspace{6pt}
\noindent\colorbox{spaet}{\parbox{0.48\textwidth}{\color{white}\bfseries\small #1}}%
\hfill%
\colorbox{english}{\parbox{0.48\textwidth}{\color{white}\bfseries\small #2}}
\vspace{3pt}
}

% Footnote style
\renewcommand{\thefootnote}{\arabic{footnote}}

\begin{document}

% Title
\begin{center}
\colorbox{spaet}{\parbox{0.48\textwidth}{\centering\color{white}\Large\bfseries\vspace{4pt}
Sequenzplan Spanisch\\SPÄTBEGINNER\\(Kl. 10/11--12/13)\\Mecklenburg-Vorpommern
\vspace{4pt}}}%
\hfill%
\colorbox{english}{\parbox{0.48\textwidth}{\centering\color{white}\Large\bfseries\vspace{4pt}
Spanish Sequence Plan\\LATE STARTERS\\(Gr. 10/11--12/13)\\Mecklenburg-Vorpommern
\vspace{4pt}}}
\end{center}

\vspace{2mm}
\begin{center}
\small Version 1.6 | Januar 2026 \hfill \textcolor{english}{Version 1.6 | January 2026}
\end{center}

\vspace{3mm}

%% DEFINITION
\secbil{1. Was ist Spätbeginner-Spanisch?}{1. What is Late Starter Spanish?}

\begin{paracol}{2}
\small
\textbf{Offizieller Begriff:} \textit{``Neu begonnene Fremdsprache''} -- es gibt keinen separaten Rahmenplan!\footnotemark[1]

\textbf{Definition:} Spanisch ab Einführungsphase (Kl. 10 bzw. Kl. 11 an Fachgymnasien/Abendgymnasien) ohne Vorkenntnisse.

\textbf{Rechtliche Grundlage:}
\begin{itemize}[noitemsep,topsep=2pt]
\item Gleicher RP wie fortgeführte FS
\item Angepasste GER-Zielniveaus (A2/B1+)
\item Eigene Bewertungsbögen (A2, B1)
\end{itemize}

\textbf{Besonderheit:} Komprimierte Progression A1→B1(+) in 2--3 Jahren.\footnotemark[2]

\switchcolumn

\small\color{english}
\textbf{Official term:} \textit{``Newly begun foreign language''} -- no separate curriculum exists!\footnotemark[1]

\textbf{Definition:} Spanish from introductory phase (Gr. 10 or 11 at vocational/evening schools) without prior knowledge.

\textbf{Legal basis:}
\begin{itemize}[noitemsep,topsep=2pt]
\item Same curriculum as continued FL
\item Adapted CEFR targets (A2/B1+)
\item Own assessment rubrics (A2, B1)
\end{itemize}

\textbf{Key feature:} Compressed A1→B1(+) in 2--3 years.\footnotemark[2]

\end{paracol}

\footnotetext[1]{\textbf{Wichtig:} Es existiert \textit{kein separater Rahmenplan} für Spätbeginner. Der Begriff ``Spätbeginner'' ist eine informelle Bezeichnung. Die offiziellen Dokumente verwenden: \textit{``In der Einführungsphase neu begonnene Fremdsprachen''}. Der gleiche \textbf{RP} = Rahmenplan Spanisch Qualifikationsphase M-V (2019) gilt für alle Kurse.}
\footnotetext[2]{\textbf{HR-SP} = Handreichung komplexe Leistungsermittlung Sprechen (2020), S. 3: ``In der Einführungsphase neu begonnene Fremdsprachen orientieren sich am Ende der Einführungsphase am Niveau A2 und am Ende der Qualifikationsphase am Zielniveau B1(+).'' \textbf{GER/GeR} = Gemeinsamer Europäischer Referenzrahmen für Sprachen (engl. CEFR).}

%% STRUCTURE
\secbil{2. Struktur, Stundentafel und GER-Progression}{2. Structure, Weekly Hours and CEFR Progression}

\vspace{2mm}

\begin{paracol}{2}
\scriptsize
\textbf{Stundentafel laut APVO M-V:}\footnotemark[3]
\begin{itemize}[noitemsep,topsep=2pt]
\item \textbf{Neu begonnene FS: 4 WStd} (vierstündig!)
\item Fortgeführte FS (GK): 3 WStd
\item Leistungskurs: 5 WStd
\end{itemize}
\textit{Spätbeginner = 4 WStd $\times$ 40 Wochen $\times$ 3 Jahre $\approx$ \textbf{480h}}
\switchcolumn
\scriptsize\color{english}
\textbf{Weekly hours per APVO M-V:}\footnotemark[3]
\begin{itemize}[noitemsep,topsep=2pt]
\item \textbf{Newly begun FL: 4 hrs/week}
\item Continued FL (basic): 3 hrs/week
\item Advanced course: 5 hrs/week
\end{itemize}
\textit{Late starters = 4 hrs $\times$ 40 weeks $\times$ 3 years $\approx$ \textbf{480h}}
\end{paracol}

\vspace{2mm}

\begin{center}
\small
\begin{tabular}{|l|c|c|c|c|c|}
\hline
\rowcolor{lightgray}
& \multicolumn{2}{c|}{\textbf{Fortgeführt / Continued}} & \multicolumn{3}{c|}{\textbf{Spätbeginner / Late Starter}} \\
\hline
\rowcolor{lightgray}
\textbf{Phase} & \textbf{GER} & \textbf{WStd} & \textbf{GER} & \textbf{WStd} & \textbf{h/Jahr} \\
\hline
Einführung / Intro & B1 & 3/5 & \textcolor{spaet}{\textbf{A2}} & \textcolor{spaet}{\textbf{4}} & $\approx$160 \\
\hline
Qualifikation 1 & B1+ & 3/5 & \textcolor{spaet}{\textbf{A2+}} & \textcolor{spaet}{\textbf{4}} & $\approx$160 \\
\hline
Qualifikation 2 & B2 & 3/5 & \textcolor{spaet}{\textbf{B1(+)}} & \textcolor{spaet}{\textbf{4}} & $\approx$160 \\
\hline
\textbf{Abitur} & \textbf{B2} & & \textcolor{spaet}{\textbf{B1(+)}} & & $\Sigma\approx$\textbf{480} \\
\hline
\end{tabular}
\end{center}

\vspace{2mm}

\begin{paracol}{2}
\small
\textbf{GER-Progression Spätbeginner:}
\begin{center}
\begin{tikzpicture}[scale=0.8]
\node[draw, fill=lightspaet, minimum width=1.2cm] at (0,0) {A1};
\draw[->, thick] (0.8,0) -- (1.5,0);
\node[draw, fill=lightspaet, minimum width=1.2cm] at (2.5,0) {A2};
\draw[->, thick] (3.3,0) -- (4,0);
\node[draw, fill=lightspaet, minimum width=1.2cm] at (5,0) {A2+};
\draw[->, thick] (5.8,0) -- (6.5,0);
\node[draw, fill=spaet, text=white, minimum width=1.4cm] at (7.7,0) {B1(+)};
\node[below] at (0,-0.5) {\tiny Start};
\node[below] at (2.5,-0.5) {\tiny Einf.};
\node[below] at (5,-0.5) {\tiny Q1};
\node[below] at (7.7,-0.5) {\tiny Abitur};
\end{tikzpicture}
\end{center}

\switchcolumn

\small\color{english}
\textbf{CEFR Progression Late Starters:}\footnotemark[3]
\begin{center}
\begin{tikzpicture}[scale=0.8]
\node[draw, fill=lightspaet, minimum width=1.2cm] at (0,0) {A1};
\draw[->, thick] (0.8,0) -- (1.5,0);
\node[draw, fill=lightspaet, minimum width=1.2cm] at (2.5,0) {A2};
\draw[->, thick] (3.3,0) -- (4,0);
\node[draw, fill=lightspaet, minimum width=1.2cm] at (5,0) {A2+};
\draw[->, thick] (5.8,0) -- (6.5,0);
\node[draw, fill=spaet, text=white, minimum width=1.4cm] at (7.7,0) {B1(+)};
\node[below] at (0,-0.5) {\tiny Start};
\node[below] at (2.5,-0.5) {\tiny Intro};
\node[below] at (5,-0.5) {\tiny Q1};
\node[below] at (7.7,-0.5) {\tiny Abitur};
\end{tikzpicture}
\end{center}

\end{paracol}

\footnotetext[3]{\textbf{APVO M-V} = Oberstufen- und Abiturprüfungsverordnung M-V (2019): ``Eine neu beginnende Fremdsprache wird \textit{vierstündig} unterrichtet.'' Fortgeführte FS (GK) = 3 WStd, LK = 5 WStd. \textbf{HR-PP} = Handreichung Paarprüfung Abitur (2025), S. 3: Zielniveau B1(+) für Spätbeginner.}

%% THEME FIELDS
\secbil{3. Die 4 Themenfelder (angepasst)}{3. The 4 Theme Fields (adapted)}

\begin{paracol}{2}
\small
Die gleichen 4 Themenfelder wie für fortgeführte Fremdsprachen gelten auch für Spätbeginner.\footnotemark[4] \textbf{Anpassungen:}
\begin{itemize}[noitemsep,topsep=2pt]
\item Reduzierte Tiefe und Komplexität
\item Einfachere Texte (A2/B1-Niveau)
\item Weniger Unterrichtsstunden verfügbar
\item Meist nur Grundkurs (kein LK)
\end{itemize}

\switchcolumn

\small\color{english}
The same 4 theme fields as for continued languages apply to late starters.\footnotemark[4] \textbf{Adaptations:}
\begin{itemize}[noitemsep,topsep=2pt]
\item Reduced depth and complexity
\item Simpler texts (A2/B1 level)
\item Fewer teaching hours available
\item Usually basic course only (no advanced)
\end{itemize}

\end{paracol}

\footnotetext[4]{RP Spanisch Sek II (2019), S. 16--18: Die vier Themenfelder gelten verbindlich für alle Kurse der Qualifikationsphase. \textbf{GK} = Grundkurs (3 WStd), \textbf{LK} = Leistungskurs (5 WStd). \textbf{WStd} = Wochenstunden.}

\vspace{2mm}

\begin{center}
\scriptsize
\begin{tabular}{|c|p{5cm}|c|c|}
\hline
\rowcolor{lightgray}
\textbf{Nr} & \textbf{Themenfeld / Theme Field} & \textbf{Fortgef.} & \textbf{Spätbeg.} \\
\hline
1 & El individuo y la sociedad \newline \textcolor{english}{\footnotesize Individual and Society} & 45h & \textcolor{spaet}{30--35h} \\
\hline
2 & La identidad nacional y la diversidad cultural \newline \textcolor{english}{\footnotesize National Identity and Cultural Diversity} & 45h & \textcolor{spaet}{30--35h} \\
\hline
3 & Aspectos actuales de la política y la sociedad \newline \textcolor{english}{\footnotesize Current Aspects of Politics and Society} & 30h & \textcolor{spaet}{20--25h} \\
\hline
4 & Desafíos modernos \newline \textcolor{english}{\footnotesize Modern Challenges} & 30h & \textcolor{spaet}{20--25h} \\
\hline
\rowcolor{lightspaet}
& \textbf{GESAMT / TOTAL} & \textbf{150h} & \textcolor{spaet}{\textbf{100--120h}} \\
\hline
\end{tabular}
\end{center}

\newpage

%% YEAR 1
\secbil{4. Jahr 1: Einführungsphase (Kl. 10/11)}{4. Year 1: Introductory Phase (Gr. 10/11)}

\begin{paracol}{2}
\small
\textbf{Zielniveau:} A2 (Ende Jahr 1)\footnotemark[5]
\par\vspace{1mm}
\textbf{Fokus:} Grundlagen aufbauen, kommunikative Basiskompetenz entwickeln

\textbf{Halbjahr 1.1 (Aug--Jan):}
\begin{itemize}[noitemsep,topsep=2pt]
\item Wo. 1--4: Grundlagen (Aussprache, Alphabet, Basisgrammatik)
\item Wo. 5--8: Yo y mi entorno (sich vorstellen, Familie)
\item Wo. 9--12: La vida cotidiana (Alltag, Uhrzeit, Aktivitäten)
\item Wo. 13--16: Mi colegio y mis amigos + Klausur
\end{itemize}

\textbf{Halbjahr 1.2 (Feb--Jul):}
\begin{itemize}[noitemsep,topsep=2pt]
\item Wo. 1--4: España: geografía básica
\item Wo. 5--8: Tradiciones y fiestas
\item Wo. 9--12: De compras y en el restaurante
\item Wo. 13--16: \textbf{Sprechprüfung} (A2) + Wiederholung
\end{itemize}

\switchcolumn

\small\color{english}
\textbf{Target level:} A2 (end of Year 1)\footnotemark[5]
\par\vspace{1mm}
\textbf{Focus:} Build foundations, develop basic communicative competence

\textbf{Semester 1.1 (Aug--Jan):}
\begin{itemize}[noitemsep,topsep=2pt]
\item Wk. 1--4: Basics (pronunciation, alphabet, basic grammar)
\item Wk. 5--8: Me and my surroundings (introduction, family)
\item Wk. 9--12: Daily life (routines, time, activities)
\item Wk. 13--16: My school and friends + Written exam
\end{itemize}

\textbf{Semester 1.2 (Feb--Jul):}
\begin{itemize}[noitemsep,topsep=2pt]
\item Wk. 1--4: Spain: basic geography
\item Wk. 5--8: Traditions and festivals
\item Wk. 9--12: Shopping and at the restaurant
\item Wk. 13--16: \textbf{Oral exam} (A2) + Review
\end{itemize}

\end{paracol}

\footnotetext[5]{Handreichung Sprechprüfung (2020), S. 3}

\vspace{2mm}

\begin{paracol}{2}
\small
\textbf{Grammatik Jahr 1:}
\begin{itemize}[noitemsep,topsep=2pt]
\item Presente de indicativo (regelmäßig + wichtige unregelmäßige)
\item Artículos, género, número
\item Possessivbegleiter
\item Pretérito perfecto (Einführung)
\item Futuro próximo (ir a + infinitivo)
\item Grundlegende Präpositionen
\end{itemize}

\switchcolumn

\small\color{english}
\textbf{Grammar Year 1:}
\begin{itemize}[noitemsep,topsep=2pt]
\item Present indicative (regular + key irregular)
\item Articles, gender, number
\item Possessive adjectives
\item Present perfect (introduction)
\item Near future (ir a + infinitive)
\item Basic prepositions
\end{itemize}

\end{paracol}

%% YEAR 2
\secbil{5. Jahr 2: Qualifikationsphase 1 (Kl. 11/12)}{5. Year 2: Qualification Phase 1 (Gr. 11/12)}

\begin{paracol}{2}
\small
\textbf{Zielniveau:} A2+ $\rightarrow$ B1 (Progression)\footnotemark[6]
\par\vspace{1mm}
\textbf{Themenfelder:} 1 + 2 (vereinfacht)

\textbf{Halbjahr 2.1: El individuo y la sociedad}
\begin{itemize}[noitemsep,topsep=2pt]
\item Wo. 1--4: Jóvenes y su mundo
\item Wo. 5--8: Relaciones y emociones
\item Wo. 9--12: Perspectivas futuras (Beruf, Studium)
\item Wo. 13--16: Klausur + Wiederholung
\end{itemize}

\textbf{Halbjahr 2.2: La identidad nacional}
\begin{itemize}[noitemsep,topsep=2pt]
\item Wo. 1--4: España: regiones y cultura
\item Wo. 5--8: Latinoamérica: introducción
\item Wo. 9--12: Diversidad cultural
\item Wo. 13--16: \textbf{Sprechprüfung} (B1)
\end{itemize}

\switchcolumn

\small\color{english}
\textbf{Target level:} A2+ $\rightarrow$ B1 (progression)\footnotemark[6]
\par\vspace{1mm}
\textbf{Theme fields:} 1 + 2 (simplified)

\textbf{Semester 2.1: Individual and Society}
\begin{itemize}[noitemsep,topsep=2pt]
\item Wk. 1--4: Young people and their world
\item Wk. 5--8: Relationships and emotions
\item Wk. 9--12: Future perspectives (career, studies)
\item Wk. 13--16: Written exam + Review
\end{itemize}

\textbf{Semester 2.2: National Identity}
\begin{itemize}[noitemsep,topsep=2pt]
\item Wk. 1--4: Spain: regions and culture
\item Wk. 5--8: Latin America: introduction
\item Wk. 9--12: Cultural diversity
\item Wk. 13--16: \textbf{Oral exam} (B1)
\end{itemize}

\end{paracol}

\footnotetext[6]{RP Spanisch Sek II (2019), S. 7: ``Am Ende der gymnasialen Oberstufe wird [...] das Niveau B2 des GeR erwartet'' -- für Spätbeginner entsprechend B1(+). \textbf{Hinweis:} Die Wochensequenzen sind eine \textit{didaktische Synthese} -- der RP gibt keine verbindliche Reihenfolge vor.}

\vspace{2mm}

\begin{paracol}{2}
\small
\textbf{Grammatik Jahr 2:}
\begin{itemize}[noitemsep,topsep=2pt]
\item Pretérito indefinido (vollständig)
\item Pretérito imperfecto
\item Kontrast indefinido/imperfecto
\item Subjuntivo presente (Einführung)
\item Condicional simple
\item Relativsätze (que, donde)
\end{itemize}

\switchcolumn

\small\color{english}
\textbf{Grammar Year 2:}
\begin{itemize}[noitemsep,topsep=2pt]
\item Preterite (complete)
\item Imperfect
\item Preterite/imperfect contrast
\item Present subjunctive (introduction)
\item Simple conditional
\item Relative clauses (que, donde)
\end{itemize}

\end{paracol}

%% YEAR 3
\secbil{6. Jahr 3: Qualifikationsphase 2 + Abitur (Kl. 12/13)}{6. Year 3: Qualification Phase 2 + Abitur (Gr. 12/13)}

\begin{paracol}{2}
\small
\textbf{Zielniveau:} B1(+) (Abitur)\footnotemark[7]
\par\vspace{1mm}
\textbf{Themenfelder:} 3 + 4 + Abitur-Vorbereitung

\textbf{Halbjahr 3.1: Aspectos actuales}
\begin{itemize}[noitemsep,topsep=2pt]
\item Wo. 1--4: Migración y sus causas
\item Wo. 5--8: Globalización
\item Wo. 9--12: Klausur (abiturähnlich)
\item Wo. 13--16: Wiederholung Thema 1+2
\end{itemize}

\textbf{Halbjahr 3.2: Desafíos + Abitur}
\begin{itemize}[noitemsep,topsep=2pt]
\item Wo. 1--4: Medio ambiente
\item Wo. 5--6: Tecnología y redes sociales
\item Wo. 7--10: \textbf{Abitur-Vorbereitung} (alle Themen)
\item Wo. 11+: Abiturprüfung
\end{itemize}

\switchcolumn

\small\color{english}
\textbf{Target level:} B1(+) (Abitur)\footnotemark[7]
\par\vspace{1mm}
\textbf{Theme fields:} 3 + 4 + Abitur preparation

\textbf{Semester 3.1: Current Aspects}
\begin{itemize}[noitemsep,topsep=2pt]
\item Wk. 1--4: Migration and its causes
\item Wk. 5--8: Globalization
\item Wk. 9--12: Written exam (Abitur-style)
\item Wk. 13--16: Review themes 1+2
\end{itemize}

\textbf{Semester 3.2: Challenges + Abitur}
\begin{itemize}[noitemsep,topsep=2pt]
\item Wk. 1--4: Environment
\item Wk. 5--6: Technology and social media
\item Wk. 7--10: \textbf{Abitur preparation} (all themes)
\item Wk. 11+: Abitur examination
\end{itemize}

\end{paracol}

\footnotetext[7]{Handreichung Paarprüfung Abitur (2025), S. 3}

\newpage

%% LITERATURE
\secbil{7. Literaturempfehlungen für Spätbeginner}{7. Literature Recommendations for Late Starters}

\begin{paracol}{2}
\small
\textbf{Kriterien für Spätbeginner-Literatur:}
\begin{itemize}[noitemsep,topsep=2pt]
\item Sprachniveau A2--B1
\item Kürzere Texte / Auszüge
\item Annotierte Ausgaben bevorzugt
\item Klare Handlung, wenig implizite Ebenen
\end{itemize}

\switchcolumn

\small\color{english}
\textbf{Criteria for late starter literature:}
\begin{itemize}[noitemsep,topsep=2pt]
\item Language level A2--B1
\item Shorter texts / excerpts
\item Annotated editions preferred
\item Clear plot, few implicit layers
\end{itemize}

\end{paracol}

\vspace{2mm}

\begin{center}
\scriptsize
\begin{tabular}{|p{3.2cm}|p{2.5cm}|p{1.5cm}|c|p{3.5cm}|}
\hline
\rowcolor{lightgray}
\textbf{Titel / Title} & \textbf{Autor} & \textbf{Niveau} & \textbf{Thema} & \textbf{Eignung / Suitability} \\
\hline
\textit{Abdel} & Enrique Páez & A2/B1 & 3 & \textcolor{spaet}{$\bigstar$ Ideal für Spätbeginner} \\
\hline
\textit{La mirada} & Juan Madrid & A2+ & 1 & Kurz, spannend, überschaubar \\
\hline
\textit{Desconexión total} & Lourdes Miquel & A2/B1 & 1 & Gruppenarbeit möglich \\
\hline
\textit{547 amigos} & Lorenzo Silva & B1 & 4 & Aktuelles Thema, kurz \\
\hline
\textit{El desertador} & J.M. Merino & B1 & 2 & Kurzgeschichte, Geschichte \\
\hline
\textit{La Casa} (Auszüge) & Paco Roca & A2+ & 1 & Novela Gráfica, visuell \\
\hline
\end{tabular}
\end{center}

\begin{paracol}{2}
\small
\textbf{Empfehlung:}\footnotemark[8] \textit{Abdel} von Enrique Páez ist besonders geeignet: Sprachlich zugänglich (A2/B1), chronologische Erzählung, 16-jähriger Protagonist ermöglicht Identifikation, Thema Migration hochaktuell.

\switchcolumn

\small\color{english}
\textbf{Recommendation:}\footnotemark[8] \textit{Abdel} by Enrique Páez is particularly suitable: Linguistically accessible (A2/B1), chronological narrative, 16-year-old protagonist enables identification, migration topic highly relevant.

\end{paracol}

\footnotetext[8]{Kommentierte Literaturempfehlungen Spanisch M-V, S. 3}

%% EXAMS
\secbil{8. Prüfungen für Spätbeginner}{8. Examinations for Late Starters}

\vspace{2mm}

\begin{center}
\small
\begin{tabular}{|l|l|l|l|}
\hline
\rowcolor{lightgray}
\textbf{Prüfung / Exam} & \textbf{Format} & \textbf{Niveau} & \textbf{Quelle / Source} \\
\hline
Sprechprüfung Einf. & Tandem 20 Min. & \textcolor{spaet}{\textbf{A2}} & HR-SP$^5$ \\
\hline
Sprechprüfung Qual. & Tandem 20 Min. & \textcolor{spaet}{\textbf{B1}} & HR-SP$^5$ \\
\hline
Abitur schriftlich & 180 Min. (GK) & \textcolor{spaet}{\textbf{B1(+)}} & APVO M-V \\
\hline
Abitur mündlich & Paarprüfung 30 Min. & \textcolor{spaet}{\textbf{B1}} & HR-PP$^7$ \\
\hline
\end{tabular}
\end{center}

\vspace{2mm}

\begin{paracol}{2}
\small
\textbf{Sprechprüfung -- 3 Teile:}\footnotemark[9]
\begin{enumerate}[noitemsep,topsep=2pt]
\item \textbf{Interview} (2 Min.) -- Aufwärmen, keine Wertung
\item \textbf{Monolog} (4 Min./Prüfling) -- Impulsgesteuert
\item \textbf{Dialog} (5 Min./Prüfling) -- Situative Diskussion
\end{enumerate}

\textbf{Bewertung:} 60\% sprachlich, 40\% inhaltlich.

\textbf{Bewertungsbögen:} Für Spätbeginner: A2 (Einführung) und B1 (Qualifikation/Abitur).

\switchcolumn

\small\color{english}
\textbf{Oral exam -- 3 parts:}\footnotemark[9]
\begin{enumerate}[noitemsep,topsep=2pt]
\item \textbf{Interview} (2 min) -- Warm-up, not graded
\item \textbf{Monologue} (4 min/student) -- Stimulus-based
\item \textbf{Dialogue} (5 min/student) -- Situational discussion
\end{enumerate}

\textbf{Grading:} 60\% linguistic, 40\% content.

\textbf{Rubrics:} For late starters: A2 (introduction) and B1 (qualification/Abitur).

\end{paracol}

\footnotetext[9]{HR-SP (2020), Anlage II: ``Bewertungsbögen A2 oder B1 oder B2''; HR-PP (2025), S. 6: Prüfungsformat mit 3 Teilen. \textbf{APVO} = Abiturprüfungsverordnung M-V.}

%% GRAMMAR OVERVIEW
\secbil{9. Grammatik-Übersicht komprimiert}{9. Grammar Overview (compressed)}

\vspace{2mm}

\begin{paracol}{2}
\scriptsize
\textbf{Hinweis:}\footnotemark[10] Der RP gibt \textit{keine} spezifischen Grammatikinhalte vor, nur das Zielniveau (``Niveau B2 des GeR''). Die folgende Progression ist eine \textbf{didaktische Synthese}.

\switchcolumn

\scriptsize\color{english}
\textbf{Note:}\footnotemark[10] The curriculum does \textit{not} specify grammar content, only the target level (``CEFR B2''). The following progression is a \textbf{didactic synthesis}.

\end{paracol}

\vspace{2mm}

\begin{center}
\scriptsize
\begin{tabular}{|l|p{5cm}|p{5cm}|}
\hline
\rowcolor{lightgray}
\textbf{Jahr} & \textbf{Deutsch (SYNTHESE)} & \textbf{English (SYNTHESIS)} \\
\hline
\textbf{1} & Presente, Perfecto, Futuro próximo, Artículo, Posesivos, Präpositionen & Present, Perfect, Near future, Articles, Possessives, Prepositions \\
\hline
\textbf{2} & Indefinido, Imperfecto, Kontrast, Subjuntivo (Einf.), Condicional, Relativsätze & Preterite, Imperfect, Contrast, Subjunctive (intro), Conditional, Relative clauses \\
\hline
\textbf{3} & Subjuntivo (Vertief.), Pluscuamperfecto, Passiv, Indirekte Rede, Konnektoren & Subjunctive (adv.), Pluperfect, Passive, Reported speech, Connectors \\
\hline
\end{tabular}
\end{center}

\footnotetext[10]{RP Spanisch Sek II (2019), S. 12 u. 21: ``ein gefestigtes Repertoire der grundlegenden grammatischen Strukturen'' und ``gemäß Niveau B2 des GeR'' -- keine Detailvorgaben.}

\newpage

%% APPENDIX: COMPARISON
\begin{center}
\colorbox{darkgray}{\parbox{0.95\textwidth}{\centering\color{white}\Large\bfseries\vspace{4pt}
ANHANG: Vergleich Regulär vs. Spätbeginner / APPENDIX: Comparison Regular vs. Late Starter
\vspace{4pt}}}
\end{center}

\vspace{4mm}

\secbil{A. Struktureller Vergleich}{A. Structural Comparison}

\vspace{2mm}

\begin{center}
\small
\begin{tabular}{|p{3.5cm}|p{4.5cm}|p{4.5cm}|}
\hline
\rowcolor{lightgray}
\textbf{Aspekt / Aspect} & \textbf{Fortgeführt / Continued} & \textbf{Spätbeginner / Late Starter} \\
\hline
Beginn / Start & Klasse 7 & Klasse 10/11 \\
\hline
Dauer bis Abitur & 6 Jahre & 2--3 Jahre \\
\hline
Gesamtstunden bis Abitur & ca. 720h (GK) / 1200h (LK) & ca. 360--480h (nur GK) \\
\hline
Einführungsphase Niveau & B1 & A2 \\
\hline
Abitur Niveau & B2 & B1(+) \\
\hline
Kurstypen & GK und LK möglich & Meist nur GK \\
\hline
Literatur-Komplexität & B1--B2 Texte, auch komplex & A2--B1 Texte, vereinfacht \\
\hline
\end{tabular}
\end{center}

\vspace{4mm}

\secbil{B. Inhaltlicher Vergleich}{B. Content Comparison}

\begin{paracol}{2}
\small
\textbf{Fortgeführte Fremdsprache:}
\begin{itemize}[noitemsep,topsep=2pt]
\item Alle 4 Themenfelder vollständig
\item Tiefgehende literarische Analyse
\item Komplexe historische Kontexte
\item Differenzierte Argumentation
\item Varietäten des Spanischen
\item \textit{El método Grönholm} (LK)
\end{itemize}

\textbf{Prüfungsanforderungen:}
\begin{itemize}[noitemsep,topsep=2pt]
\item Abitur: 270 Min. (LK) / 180 Min. (GK)
\item AFB-Verteilung: 20/45/35
\item Komplexe Textanalyse erwartet
\item Stilmittel-Analyse
\end{itemize}

\switchcolumn

\small\color{english}
\textbf{Continued Foreign Language:}
\begin{itemize}[noitemsep,topsep=2pt]
\item All 4 theme fields complete
\item In-depth literary analysis
\item Complex historical contexts
\item Differentiated argumentation
\item Spanish language varieties
\item \textit{El método Grönholm} (advanced)
\end{itemize}

\textbf{Exam requirements:}
\begin{itemize}[noitemsep,topsep=2pt]
\item Abitur: 270 min (adv.) / 180 min (basic)
\item AFB distribution: 20/45/35
\item Complex text analysis expected
\item Stylistic device analysis
\end{itemize}

\end{paracol}

\vspace{3mm}

\begin{paracol}{2}
\small
\textbf{Spätbeginner:}
\begin{itemize}[noitemsep,topsep=2pt]
\item Alle 4 Themenfelder in reduzierter Tiefe
\item Grundlegende Textarbeit
\item Vereinfachte historische Überblicke
\item Klare, strukturierte Argumentation
\item Fokus auf Standardsprache
\item \textit{Abdel}, \textit{La mirada} (einfachere Texte)
\end{itemize}

\textbf{Prüfungsanforderungen:}
\begin{itemize}[noitemsep,topsep=2pt]
\item Abitur: 180 Min. (nur GK)
\item AFB-Verteilung: 25/50/25 (mehr AFB I)
\item Grundlegende Texterschließung
\item Weniger Stilmittel-Fokus
\end{itemize}

\switchcolumn

\small\color{english}
\textbf{Late Starters:}
\begin{itemize}[noitemsep,topsep=2pt]
\item All 4 theme fields at reduced depth
\item Basic text work
\item Simplified historical overviews
\item Clear, structured argumentation
\item Focus on standard language
\item \textit{Abdel}, \textit{La mirada} (simpler texts)
\end{itemize}

\textbf{Exam requirements:}
\begin{itemize}[noitemsep,topsep=2pt]
\item Abitur: 180 min (basic only)
\item AFB distribution: 25/50/25 (more AFB I)
\item Basic text comprehension
\item Less stylistic device focus
\end{itemize}

\end{paracol}

\vspace{4mm}

\secbil{C. Wichtige Hinweise}{C. Important Notes}

\begin{paracol}{2}
\small
\begin{enumerate}[noitemsep,topsep=2pt]
\item Spätbeginner können \textbf{nicht} als Leistungskurs gewählt werden
\item Die \textbf{gleichen Themenfelder} gelten, aber mit angepasstem Niveau
\item \textbf{Sprechprüfungen} sind auch für Spätbeginner verpflichtend
\item Die \textbf{Paarprüfung im Abitur} ist auch für Spätbeginner möglich
\item \textbf{Bewertungsbögen} müssen dem Niveau entsprechen (A2/B1)
\end{enumerate}

\switchcolumn

\small\color{english}
\begin{enumerate}[noitemsep,topsep=2pt]
\item Late starters \textbf{cannot} be chosen as advanced course
\item The \textbf{same theme fields} apply, but at adapted level
\item \textbf{Oral exams} are also mandatory for late starters
\item The \textbf{pair exam in Abitur} is also available for late starters
\item \textbf{Assessment rubrics} must match the level (A2/B1)
\end{enumerate}

\end{paracol}

\vspace{6mm}

%% ABBREVIATIONS
\vspace{4mm}
\secbil{Abkürzungsverzeichnis}{Abbreviations}

\vspace{2mm}

\begin{center}
\scriptsize
\begin{tabular}{|l|p{5cm}|p{5cm}|}
\hline
\rowcolor{lightgray}
\textbf{Abk.} & \textbf{Deutsch} & \textbf{English} \\
\hline
GER/GeR & Gemeinsamer Europäischer Referenzrahmen & Common European Framework of Reference (CEFR) \\
\hline
\textbf{APVO M-V} & \textbf{Oberstufen- und Abiturprüfungsverordnung} & \textbf{Upper Secondary \& Abitur Exam Regulations} \\
\hline
GK & Grundkurs (fortgef. FS: 3 WStd) & Basic Course (continued FL: 3 hrs/wk) \\
\hline
\textcolor{spaet}{\textbf{Neu beg. FS}} & \textcolor{spaet}{\textbf{Spätbeginner: 4 WStd}} & \textcolor{spaet}{\textbf{Late starter: 4 hrs/wk}} \\
\hline
LK & Leistungskurs (5 WStd) & Advanced Course (5 hrs/wk) \\
\hline
AFB & Anforderungsbereich (I/II/III) & Requirement Level (I/II/III) \\
\hline
RP & Rahmenplan (2019) & Curriculum Framework (2019) \\
\hline
HR-SP & Handreichung Sprechprüfung (2020) & Oral Exam Guidelines (2020) \\
\hline
HR-PP & Handreichung Paarprüfung (2025) & Pair Exam Guidelines (2025) \\
\hline
\end{tabular}
\end{center}

\newpage

%% APPENDIX: PRÜFUNGSORDNUNG
\begin{center}
\colorbox{spaet}{\parbox{0.95\textwidth}{\centering\color{white}\Large\bfseries\vspace{4pt}
ANHANG D: Prüfungsordnung im Detail / APPENDIX D: Examination Regulations in Detail
\vspace{4pt}}}
\end{center}

\vspace{4mm}

\secbil{D1. Klausuren in der Oberstufe}{D1. Written Exams in Upper Secondary}

\vspace{2mm}

\begin{paracol}{2}
\small
\textbf{Anzahl (APVO M-V):}
\begin{itemize}[noitemsep,topsep=2pt]
\item \textbf{1 Klausur pro Halbjahr} (Grundkurs)
\item Pro Jahr: max. 1 Klausur ersetzbar durch Sprechprüfung
\item Einführungsphase: 1 Ersetzung möglich
\item Qualifikationsphase: 1 Ersetzung möglich
\end{itemize}

\textbf{Dauer Klausuren:}
\begin{itemize}[noitemsep,topsep=2pt]
\item Einführungsphase: 90 Min.
\item Qualifikationsphase GK: 135 Min.
\item Qualifikationsphase LK: 180 Min.
\end{itemize}

\switchcolumn

\small\color{english}
\textbf{Number (APVO M-V):}
\begin{itemize}[noitemsep,topsep=2pt]
\item \textbf{1 written exam per semester} (basic course)
\item Per year: max. 1 exam replaceable by oral exam
\item Introductory phase: 1 replacement possible
\item Qualification phase: 1 replacement possible
\end{itemize}

\textbf{Exam duration:}
\begin{itemize}[noitemsep,topsep=2pt]
\item Introductory phase: 90 min
\item Qualification phase basic: 135 min
\item Qualification phase advanced: 180 min
\end{itemize}

\end{paracol}

\vspace{3mm}

\secbil{D2. Sprechprüfung -- Verbindliche Regelungen}{D2. Oral Exam -- Binding Regulations}

\vspace{2mm}

\begin{paracol}{2}
\small
\textbf{Rechtliche Grundlage:}\footnotemark[11]
\begin{itemize}[noitemsep,topsep=2pt]
\item §22 Abs. 5 APVO M-V: \textbf{Mindestens 1$\times$ pro SuS verpflichtend}
\item VV Sprechprüfung vom 24.07.2020
\item HR-SP: Handreichung Sprechprüfung (2020)
\end{itemize}

\textbf{Format:}
\begin{itemize}[noitemsep,topsep=2pt]
\item \textbf{Tandem-Prüfung} (2 SuS gleichzeitig)
\item Gesamtdauer: \textbf{20 Minuten}
\item 2 Prüfer (Fachprüfer + Protokollant)
\end{itemize}

\switchcolumn

\small\color{english}
\textbf{Legal basis:}\footnotemark[11]
\begin{itemize}[noitemsep,topsep=2pt]
\item §22 Para. 5 APVO M-V: \textbf{At least once per student mandatory}
\item Administrative regulation of 24.07.2020
\item HR-SP: Oral Exam Guidelines (2020)
\end{itemize}

\textbf{Format:}
\begin{itemize}[noitemsep,topsep=2pt]
\item \textbf{Tandem exam} (2 students simultaneously)
\item Total duration: \textbf{20 minutes}
\item 2 examiners (subject examiner + note-taker)
\end{itemize}

\end{paracol}

\vspace{2mm}

\begin{center}
\small
\textbf{Prüfungsablauf / Exam Structure:}
\vspace{1mm}

\begin{tabular}{|c|p{3.5cm}|c|p{6cm}|}
\hline
\rowcolor{lightgray}
\textbf{Teil} & \textbf{Name} & \textbf{Dauer} & \textbf{Beschreibung / Description} \\
\hline
1 & Interview & 2 Min. & Aufwärmen, persönliche Fragen -- \textbf{keine Bewertung!} \newline \textcolor{english}{\footnotesize Warm-up, personal questions -- \textbf{not graded!}} \\
\hline
2 & Monolog & 4 Min./SuS & Impulsgesteuerte Präsentation (Bild/Text) \newline \textcolor{english}{\footnotesize Stimulus-based presentation (image/text)} \\
\hline
3 & Dialog & 5 Min./SuS & Diskussion/Rollenspiel zu einem Thema \newline \textcolor{english}{\footnotesize Discussion/role-play on a topic} \\
\hline
\end{tabular}
\end{center}

\vspace{2mm}

\begin{paracol}{2}
\small
\textbf{Bewertung:}
\begin{itemize}[noitemsep,topsep=2pt]
\item \textbf{60\% sprachliche Leistung}
\begin{itemize}[noitemsep,topsep=1pt]
\item[-] Aussprache/Intonation
\item[-] Wortschatz/Ausdruck
\item[-] Grammatische Korrektheit
\item[-] Flüssigkeit/Interaktion
\end{itemize}
\item \textbf{40\% inhaltliche Leistung}
\begin{itemize}[noitemsep,topsep=1pt]
\item[-] Aufgabenerfüllung
\item[-] Argumentation
\item[-] Themenentwicklung
\end{itemize}
\end{itemize}

\textbf{Bewertungsbögen (verbindlich):}\footnotemark[12]
\begin{itemize}[noitemsep,topsep=2pt]
\item \textcolor{spaet}{\textbf{A2}} -- Einführungsphase Spätbeginner
\item \textcolor{spaet}{\textbf{B1}} -- Qualifikation Spätbeginner
\item B2 -- Fortgeführte Fremdsprache
\end{itemize}

\switchcolumn

\small\color{english}
\textbf{Grading:}
\begin{itemize}[noitemsep,topsep=2pt]
\item \textbf{60\% linguistic performance}
\begin{itemize}[noitemsep,topsep=1pt]
\item[-] Pronunciation/intonation
\item[-] Vocabulary/expression
\item[-] Grammatical accuracy
\item[-] Fluency/interaction
\end{itemize}
\item \textbf{40\% content performance}
\begin{itemize}[noitemsep,topsep=1pt]
\item[-] Task completion
\item[-] Argumentation
\item[-] Topic development
\end{itemize}
\end{itemize}

\textbf{Assessment rubrics (binding):}\footnotemark[12]
\begin{itemize}[noitemsep,topsep=2pt]
\item \textcolor{spaet}{\textbf{A2}} -- Introductory phase late starters
\item \textcolor{spaet}{\textbf{B1}} -- Qualification late starters
\item B2 -- Continued foreign language
\end{itemize}

\end{paracol}

\footnotetext[11]{VV vom 24.07.2020, Mittl.bl. BM M-V 2020 S. 290: ``Die Abiturprüfungsverordnung [...] schreibt in §22 Absatz 5 eine komplexe Leistungsermittlung im Kompetenzbereich Sprechen vor.''}
\footnotetext[12]{HR-SP (2020), Anlage II: Drei Bewertungsbögen (A2, B1, B2) -- das GER-Zielniveau des Kurses bestimmt den Bogen. Spätbeginner nutzen A2 (Einführung) bzw. B1 (Qualifikation).}

\vspace{3mm}

\secbil{D2b. Bewertungsbögen im Detail}{D2b. Assessment Rubrics in Detail}

\vspace{2mm}

\begin{center}
\small
\textbf{Verfügbare Bewertungsbögen / Available Assessment Rubrics:}
\vspace{1mm}

\begin{tabular}{|l|c|c|l|}
\hline
\rowcolor{lightgray}
\textbf{Prüfung / Exam} & \textbf{Bögen} & \textbf{Spätbeg.} & \textbf{Quelle / Source} \\
\hline
Sprechprüfung (Unterricht) & A2, B1, B2 & \textcolor{spaet}{\textbf{A2/B1}} & HR-SP Anlage II \\
\hline
Abitur Paarprüfung & B1, B2 & \textcolor{spaet}{\textbf{B1}} & HR-PP Anlage I \\
\hline
\end{tabular}
\end{center}

\vspace{2mm}

\begin{paracol}{2}
\small
\textbf{Struktur der Bewertungsbögen:}\footnotemark[12]
\begin{itemize}[noitemsep,topsep=2pt]
\item \textbf{1 Bogen Gesprächsleiter}
\item \textbf{2 Bögen Protokollant:}
\begin{itemize}[noitemsep,topsep=1pt]
\item[-] I. Monologisches Sprechen
\item[-] II. Dialogisches Sprechen
\end{itemize}
\item Format: Excel-Arbeitsmappen
\end{itemize}

\textbf{Bewertungskriterien:}
\begin{enumerate}[noitemsep,topsep=2pt]
\item \textbf{Inhalt/Aufgabenerfüllung (40\%)}
\item \textbf{Darstellung/Sprache (60\%):}
\begin{itemize}[noitemsep,topsep=1pt]
\item[-] Flüssigkeit \& Interaktion
\item[-] Spektrum \& Richtigkeit
\end{itemize}
\end{enumerate}

\switchcolumn

\small\color{english}
\textbf{Structure of assessment rubrics:}\footnotemark[12]
\begin{itemize}[noitemsep,topsep=2pt]
\item \textbf{1 rubric for lead examiner}
\item \textbf{2 rubrics for note-taker:}
\begin{itemize}[noitemsep,topsep=1pt]
\item[-] I. Monological speaking
\item[-] II. Dialogical speaking
\end{itemize}
\item Format: Excel workbooks
\end{itemize}

\textbf{Assessment criteria:}
\begin{enumerate}[noitemsep,topsep=2pt]
\item \textbf{Content/Task completion (40\%)}
\item \textbf{Presentation/Language (60\%):}
\begin{itemize}[noitemsep,topsep=1pt]
\item[-] Fluency \& Interaction
\item[-] Range \& Accuracy
\end{itemize}
\end{enumerate}

\end{paracol}

\vspace{2mm}

\begin{center}
\scriptsize
\textbf{Bezugsquellen / Download Sources:}\\
SIP-Datentauschportal: ``Sprechleistung\_Moderne Fremdsprachen''\\
itslearning: ``Sprechen testen in den Fremdsprachen SEK II'' (HR-SP) | ``Paarprüfung im Abitur'' (HR-PP)
\end{center}

\vspace{2mm}

\begin{paracol}{2}
\scriptsize
\textbf{Kombination Sprechen + Hörverstehen:}\\
Falls kombiniert: 80\% Sprechen, 20\% Hör-/Hörsehverstehen (§22 Abs. 5 APVO).

\switchcolumn

\scriptsize\color{english}
\textbf{Combination Speaking + Listening:}\\
If combined: 80\% Speaking, 20\% Listening (§22 Para. 5 APVO).

\end{paracol}

\vspace{3mm}

\secbil{D3. Abiturprüfung schriftlich}{D3. Written Abitur Exam}

\vspace{2mm}

\begin{paracol}{2}
\small
\textbf{Dauer nach APVO M-V:}
\begin{itemize}[noitemsep,topsep=2pt]
\item Grundkurs: \textbf{180 Minuten}
\item Leistungskurs: \textbf{270 Minuten}
\item \textcolor{spaet}{Spätbeginner: nur GK (180 Min.)}
\end{itemize}

\textbf{Aufgabenformate:}
\begin{itemize}[noitemsep,topsep=2pt]
\item Textaufgabe (Pflicht)
\item Leseverstehen
\item Hörverstehen (optional)
\item Sprachmittlung (optional)
\end{itemize}

\textbf{AFB-Verteilung:}
\begin{itemize}[noitemsep,topsep=2pt]
\item \textcolor{spaet}{Spätbeginner: 25\% / 50\% / 25\%}
\item Fortgeführt: 20\% / 45\% / 35\%
\end{itemize}

\switchcolumn

\small\color{english}
\textbf{Duration per APVO M-V:}
\begin{itemize}[noitemsep,topsep=2pt]
\item Basic course: \textbf{180 minutes}
\item Advanced course: \textbf{270 minutes}
\item \textcolor{spaet}{Late starters: basic only (180 min)}
\end{itemize}

\textbf{Task formats:}
\begin{itemize}[noitemsep,topsep=2pt]
\item Text task (mandatory)
\item Reading comprehension
\item Listening comprehension (optional)
\item Mediation (optional)
\end{itemize}

\textbf{AFB distribution:}
\begin{itemize}[noitemsep,topsep=2pt]
\item \textcolor{spaet}{Late starters: 25\% / 50\% / 25\%}
\item Continued: 20\% / 45\% / 35\%
\end{itemize}

\end{paracol}

\vspace{3mm}

\secbil{D4. Abitur Paarprüfung (mündlich)}{D4. Abitur Pair Exam (oral)}

\vspace{2mm}

\begin{paracol}{2}
\small
\textbf{Rechtliche Grundlage:}\footnotemark[13]
\begin{itemize}[noitemsep,topsep=2pt]
\item §25 Abs. 8 APVO M-V (Kombination)
\item §38 Abs. 6 APVO M-V (mündl. Prüfung)
\item VV Paarprüfung vom 03.12.2025
\item HR-PP: Handreichung Paarprüfung (2025)
\end{itemize}

\textbf{Anwendungsfälle:}
\begin{itemize}[noitemsep,topsep=2pt]
\item \textbf{P5 (mündliche Abiturprüfung):} vollständig mündlich
\item \textbf{P4 (Kombination):} verkürzte schriftliche + mündliche Prüfung
\end{itemize}

\switchcolumn

\small\color{english}
\textbf{Legal basis:}\footnotemark[13]
\begin{itemize}[noitemsep,topsep=2pt]
\item §25 Para. 8 APVO M-V (combination)
\item §38 Para. 6 APVO M-V (oral exam)
\item Administrative regulation of 03.12.2025
\item HR-PP: Pair Exam Guidelines (2025)
\end{itemize}

\textbf{Application cases:}
\begin{itemize}[noitemsep,topsep=2pt]
\item \textbf{P5 (oral Abitur exam):} fully oral
\item \textbf{P4 (combination):} shortened written + oral exam
\end{itemize}

\end{paracol}

\vspace{2mm}

\begin{center}
\small
\textbf{Prüfungsablauf Paarprüfung Abitur / Exam Structure Pair Exam:}
\vspace{1mm}

\begin{tabular}{|c|p{3.5cm}|c|p{6cm}|}
\hline
\rowcolor{lightgray}
\textbf{Teil} & \textbf{Name} & \textbf{Dauer} & \textbf{Beschreibung / Description} \\
\hline
1 & Interview & 2 Min. & Kurzer Einstieg, Auflockerung -- \textbf{keine Bewertung!} \newline \textcolor{english}{\footnotesize Brief introduction, warm-up -- \textbf{not graded!}} \\
\hline
2 & Monolog & 4 Min./SuS & Längerer Gesprächsbeitrag auf Basis eines Impulses \newline \textcolor{english}{\footnotesize Extended contribution based on stimulus} \\
\hline
3 & Dialog & 5 Min./SuS & Interaktion zwischen beiden Prüflingen \newline \textcolor{english}{\footnotesize Interaction between both examinees} \\
\hline
\multicolumn{2}{|c|}{\textbf{Gesamt / Total}} & \textbf{ca. 30 Min.} & Vorbereitungszeit: 20 Min. (beide Prüflinge gemeinsam) \\
\hline
\end{tabular}
\end{center}

\vspace{2mm}

\begin{paracol}{2}
\small
\textbf{Bewertung Paarprüfung:}
\begin{itemize}[noitemsep,topsep=2pt]
\item \textbf{Individuelle Bewertung} jedes Prüflings
\item 60\% sprachlich / 40\% inhaltlich
\item Bewertungsbögen: \textcolor{spaet}{\textbf{B1}} (Spätbeginner) oder B2
\item Keine Beeinflussung durch Partner
\end{itemize}

\textbf{Besonderheiten:}
\begin{itemize}[noitemsep,topsep=2pt]
\item Impulsmaterial von Prüfern gewählt
\item Beide Prüflinge erhalten \textbf{unterschiedliche} Teilimpulse
\item Dialog: \textbf{eigenständige Interaktion} erwartet
\item Prüfer greifen nur bei Bedarf ein
\end{itemize}

\switchcolumn

\small\color{english}
\textbf{Pair exam grading:}
\begin{itemize}[noitemsep,topsep=2pt]
\item \textbf{Individual grading} of each examinee
\item 60\% linguistic / 40\% content
\item Rubrics: \textcolor{spaet}{\textbf{B1}} (late starters) or B2
\item No influence from partner
\end{itemize}

\textbf{Special features:}
\begin{itemize}[noitemsep,topsep=2pt]
\item Stimulus material chosen by examiners
\item Both examinees receive \textbf{different} sub-stimuli
\item Dialogue: \textbf{autonomous interaction} expected
\item Examiners intervene only if necessary
\end{itemize}

\end{paracol}

\footnotetext[13]{VV vom 03.12.2025, Mittl.bl. BM M-V 2025 S. 239: ``Die Abiturprüfungsverordnung [...] ermöglicht die Umsetzung der Paarprüfung im Kompetenzbereich Sprechen im Rahmen der Abiturprüfungen.''}

\vspace{3mm}

\secbil{D5. Zusammenfassung: Prüfungen für Spätbeginner}{D5. Summary: Exams for Late Starters}

\vspace{2mm}

\begin{center}
\small
\begin{tabular}{|p{2.8cm}|c|c|c|p{4.8cm}|}
\hline
\rowcolor{lightspaet}
\textbf{Prüfung / Exam} & \textbf{Anzahl} & \textbf{Dauer} & \textbf{Niveau} & \textbf{Besonderheit} \\
\hline
Klausur Einführung (Kl. 10) & 2/Jahr & 90 Min. & A2 & 1$\times$/Jahr ersetzbar durch Sprechprüfung \\
\hline
Klausur Qualifikation (Kl. 11-12) & 2/Jahr & 135 Min. & A2+/B1 & 1$\times$/Jahr ersetzbar durch Sprechprüfung \\
\hline
\textbf{Sprechprüfung} & \textbf{$\geq$1 ges.} & 20 Min. & A2/B1 & \textbf{Pflicht!} Empfehlung: Kl. 11 oder 12 \\
\hline
Abitur schriftlich (Kl. 12) & 1 & 180 Min. & B1(+) & Nur auf grundlegendem Niveau \\
\hline
Abitur mündlich P5 & 1 & 30 Min. & B1 & Optional: Paarprüfung \\
\hline
Abitur Kombination P4 & 1 & variabel & B1 & Verkürzt schriftl. + mündl. \\
\hline
\end{tabular}
\end{center}

\vspace{2mm}

\begin{paracol}{2}
\small
\textbf{Wichtig: Spätbeginner = eigene Kategorie!}\footnotemark[15]
\begin{itemize}[noitemsep,topsep=2pt]
\item ``Neu begonnene FS'' = \textbf{4 WStd}
\item Weder GK (3 WStd) noch LK (5 WStd)!
\item Prüfungen: wie GK-Niveau behandelt
\item Klausurlänge nach \textbf{Jahrgang}, nicht Kursart
\end{itemize}

\textbf{Sprechprüfung -- Timing:}\footnotemark[16]
\begin{itemize}[noitemsep,topsep=2pt]
\item Mind. 1$\times$ pro SuS verpflichtend
\item \textbf{Empfehlung:} Kl. 11 oder 12 (nicht Kl. 10)
\item 2./3. FS: eher im 2. Schulhalbjahr
\item Grund: Korrekturentlastung Englisch-LK
\end{itemize}

\switchcolumn

\small\color{english}
\textbf{Important: Late starters = separate category!}\footnotemark[15]
\begin{itemize}[noitemsep,topsep=2pt]
\item ``Newly begun FL'' = \textbf{4 hrs/wk}
\item Neither basic (3) nor advanced (5) course!
\item Exams: treated as basic level
\item Exam duration by \textbf{grade level}, not course type
\end{itemize}

\textbf{Oral exam -- Timing:}\footnotemark[16]
\begin{itemize}[noitemsep,topsep=2pt]
\item At least 1$\times$ per student mandatory
\item \textbf{Recommendation:} Grade 11 or 12 (not 10)
\item 2nd/3rd FL: preferably in 2nd semester
\item Reason: reduce grading load for English adv.
\end{itemize}

\end{paracol}

\footnotetext[15]{APVO M-V (2019): ``Eine neu beginnende Fremdsprache wird \textit{vierstündig} unterrichtet.'' Dies ist KEINE Kursart (GK/LK), sondern eine eigene Kategorie mit 4 WStd. Für Prüfungszwecke wird das grundlegende Niveau (GK-äquivalent) angewandt.}
\footnotetext[16]{HR-SP (2020), S. 6: ``Es wird empfohlen, im zweiten Schulhalbjahr die zweiten bzw. dritten Fremdsprachen zu überprüfen.'' Grund: Englisch-Lehrkräfte korrigieren im 4. Halbjahr die schriftlichen Abiturprüfungen.}

\vspace{4mm}

\secbil{D6. Vergleich: Identisch vs. Unterschiedlich}{D6. Comparison: Identical vs. Different}

\vspace{2mm}

\begin{paracol}{2}
\small
\textbf{IDENTISCH mit fortgeführter FS:}\footnotemark[14]
\begin{itemize}[noitemsep,topsep=2pt]
\item Format Sprechprüfung (20 Min., Tandem)
\item Format Paarprüfung Abitur (30 Min.)
\item 3-Teile-Struktur (Interview/Monolog/Dialog)
\item Bewertung 60\% sprachlich / 40\% inhaltlich
\item Mind. 1$\times$ Sprechprüfung verpflichtend
\item Teil 1 (Interview) nicht bewertet
\item 2 Prüfer (Gesprächsleiter + Protokollant)
\item Klausurersetzung möglich (1$\times$/Jahr)
\end{itemize}

\switchcolumn

\small\color{english}
\textbf{IDENTICAL to continued FL:}\footnotemark[14]
\begin{itemize}[noitemsep,topsep=2pt]
\item Oral exam format (20 min, tandem)
\item Abitur pair exam format (30 min)
\item 3-part structure (Interview/Monologue/Dialogue)
\item Grading 60\% linguistic / 40\% content
\item At least 1$\times$ oral exam mandatory
\item Part 1 (Interview) not graded
\item 2 examiners (lead + note-taker)
\item Exam replacement possible (1$\times$/year)
\end{itemize}

\end{paracol}

\vspace{2mm}

\begin{center}
\small
\textbf{UNTERSCHIEDE / DIFFERENCES:}
\vspace{1mm}

\begin{tabular}{|p{4cm}|p{4cm}|p{4cm}|}
\hline
\rowcolor{lightgray}
\textbf{Aspekt / Aspect} & \textbf{Fortgeführt / Continued} & \textbf{\textcolor{spaet}{Spätbeginner / Late Starter}} \\
\hline
GER Ende Einführung & B1 & \textcolor{spaet}{\textbf{A2}} \\
\hline
GER Ende Qualifikation & B2 & \textcolor{spaet}{\textbf{B1(+)}} \\
\hline
Bewertungsbogen Sprechpr. & B1 / B2 & \textcolor{spaet}{\textbf{A2 / B1}} \\
\hline
Bewertungsbogen Abitur & B2 & \textcolor{spaet}{\textbf{B1}} \\
\hline
Leistungskurs möglich? & Ja & \textcolor{spaet}{\textbf{Nein}} \\
\hline
Abitur schriftl. max. & 270 Min. (LK) & \textcolor{spaet}{\textbf{180 Min. (nur GK)}} \\
\hline
\end{tabular}
\end{center}

\vspace{2mm}

\begin{paracol}{2}
\scriptsize
\textbf{Fazit:} Die Prüfungs\textit{formate} und \textit{-verfahren} sind \textbf{identisch}. Nur die \textbf{GER-Zielniveaus} und damit die \textbf{Bewertungsbögen} unterscheiden sich. Spätbeginner nutzen A2 (Einführung) bzw. B1 (Qualifikation/Abitur).

\switchcolumn

\scriptsize\color{english}
\textbf{Conclusion:} The exam \textit{formats} and \textit{procedures} are \textbf{identical}. Only the \textbf{CEFR target levels} and thus the \textbf{assessment rubrics} differ. Late starters use A2 (introduction) or B1 (qualification/Abitur).

\end{paracol}

\footnotetext[14]{HR-SP (2020) und HR-PP (2025): Prüfungsformat und -ablauf gelten für alle modernen Fremdsprachen gleichermaßen. Einziger Unterschied: Die Bewertungsbögen werden nach GER-Zielniveau gewählt (A2, B1 oder B2).}

\vspace{4mm}

\secbil{D7. Praxishinweise: Jahrgangsübergreifende Kurse}{D7. Practical Notes: Mixed-Grade Courses}

\vspace{2mm}

\begin{paracol}{2}
\small
\textbf{Situation an kleineren Schulen:}\footnotemark[17]

An Gesamtschulen werden Spätbeginner-Kurse oft \textbf{jahrgangsübergreifend} unterrichtet:
\begin{itemize}[noitemsep,topsep=2pt]
\item Kl. 10, 11, 12 \textbf{gemeinsam} in einem Kurs
\item Alle mit \textbf{4 WStd} (APVO-konform)
\item Gleicher Unterricht, aber unterschiedliche Prüfungsanforderungen
\end{itemize}

\textbf{Ist das zulässig?}

\textbf{Ja}, wenn folgende Bedingungen erfüllt sind:
\begin{enumerate}[noitemsep,topsep=2pt]
\item Alle SuS erhalten \textbf{4 WStd}
\item \textbf{Klausuren differenziert} nach Jahrgang:
\begin{itemize}[noitemsep,topsep=1pt]
\item[-] Kl. 10: 90 Min. (A2-Niveau)
\item[-] Kl. 11-12: 135 Min. (A2+/B1)
\end{itemize}
\item Bewertung nach \textbf{individuellem Stand}
\item Mind. 1 Sprechprüfung pro SuS
\end{enumerate}

\switchcolumn

\small\color{english}
\textbf{Situation at smaller schools:}\footnotemark[17]

At comprehensive schools, late starter courses are often taught \textbf{across grade levels}:
\begin{itemize}[noitemsep,topsep=2pt]
\item Grades 10, 11, 12 \textbf{together} in one course
\item All with \textbf{4 hrs/wk} (APVO-compliant)
\item Same instruction, but different exam requirements
\end{itemize}

\textbf{Is this permissible?}

\textbf{Yes}, if the following conditions are met:
\begin{enumerate}[noitemsep,topsep=2pt]
\item All students receive \textbf{4 hrs/wk}
\item \textbf{Exams differentiated} by grade:
\begin{itemize}[noitemsep,topsep=1pt]
\item[-] Gr. 10: 90 min (A2 level)
\item[-] Gr. 11-12: 135 min (A2+/B1)
\end{itemize}
\item Grading by \textbf{individual level}
\item At least 1 oral exam per student
\end{enumerate}

\end{paracol}

\vspace{2mm}

\begin{center}
\small
\textbf{Prüfungsgestaltung bei gemischten Kursen / Exam Design for Mixed Courses:}
\vspace{1mm}

\begin{tabular}{|p{2.5cm}|p{4.5cm}|p{4.5cm}|}
\hline
\rowcolor{lightgray}
\textbf{Jahrgang} & \textbf{Klausur} & \textbf{Sprechprüfung} \\
\hline
Kl. 10 (1. Jahr) & 90 Min., A2-Bogen, 2/Jahr & Optional (Einführung) \\
\hline
Kl. 11 (2. Jahr) & 135 Min., A2+/B1-Bogen, 2/Jahr & \textbf{Empfohlen} (ersetzt 1 Klausur) \\
\hline
Kl. 12 (3. Jahr) & 135 Min., B1-Bogen, 2/Jahr & Spätestens jetzt \textbf{Pflicht} \\
\hline
\end{tabular}
\end{center}

\vspace{2mm}

\begin{paracol}{2}
\scriptsize
\textbf{Praxis-Tipps:}
\begin{itemize}[noitemsep,topsep=2pt]
\item Klausuren: \textbf{gleiche Themen}, aber unterschiedliche Länge/Komplexität
\item Sprechprüfung: Tandems \textbf{aus gleichem Jahrgang} bilden
\item Oder: jahrgangsgemischt mit angepasstem Bewertungsbogen
\item Dokumentation: Jahrgang und GER-Niveau pro SuS festhalten
\end{itemize}

\textbf{Falls kein Abitur möglich:}\\
Auch ohne prüfungsbefugte Lehrkraft: Kurse können durchgeführt werden. Die Sprechprüfung kann ggf. extern abgenommen werden (Kooperation mit anderer Schule).

\switchcolumn

\scriptsize\color{english}
\textbf{Practical tips:}
\begin{itemize}[noitemsep,topsep=2pt]
\item Exams: \textbf{same topics}, but different length/complexity
\item Oral exam: form tandems \textbf{from same grade}
\item Or: mixed grades with adapted rubric
\item Documentation: record grade and CEFR level per student
\end{itemize}

\textbf{If no Abitur possible:}\\
Even without certified examiner: courses can be conducted. Oral exams may be administered externally (cooperation with another school).

\end{paracol}

\footnotetext[17]{Jahrgangsübergreifender Unterricht ist in M-V zulässig, wenn die APVO-Vorgaben (4 WStd, Prüfungsformate) eingehalten werden. Die Schulkonferenz muss dies beschließen. Pragmatisch: Gemeinsamer Unterricht, differenzierte Prüfungen.}

\vspace{6mm}

%% ABBREVIATIONS
\vspace{4mm}
\secbil{Abkürzungsverzeichnis}{Abbreviations}

\vspace{2mm}

\begin{center}
\scriptsize
\begin{tabular}{|l|p{5cm}|p{5cm}|}
\hline
\rowcolor{lightgray}
\textbf{Abk.} & \textbf{Deutsch} & \textbf{English} \\
\hline
GER/GeR & Gemeinsamer Europäischer Referenzrahmen & Common European Framework of Reference (CEFR) \\
\hline
\textbf{APVO M-V} & \textbf{Oberstufen- und Abiturprüfungsverordnung} & \textbf{Upper Secondary \& Abitur Exam Regulations} \\
\hline
GK & Grundkurs (fortgef. FS: 3 WStd) & Basic Course (continued FL: 3 hrs/wk) \\
\hline
\textcolor{spaet}{\textbf{Neu beg. FS}} & \textcolor{spaet}{\textbf{Spätbeginner: 4 WStd}} & \textcolor{spaet}{\textbf{Late starter: 4 hrs/wk}} \\
\hline
LK & Leistungskurs (5 WStd) & Advanced Course (5 hrs/wk) \\
\hline
AFB & Anforderungsbereich (I/II/III) & Requirement Level (I/II/III) \\
\hline
RP & Rahmenplan (2019) & Curriculum Framework (2019) \\
\hline
HR-SP & Handreichung Sprechprüfung (2020) & Oral Exam Guidelines (2020) \\
\hline
HR-PP & Handreichung Paarprüfung (2025) & Pair Exam Guidelines (2025) \\
\hline
VV & Verwaltungsvorschrift & Administrative Regulation \\
\hline
P4/P5 & Prüfungsfach 4/5 (Abitur) & Exam subject 4/5 (Abitur) \\
\hline
SuS & Schülerinnen und Schüler & Students \\
\hline
\end{tabular}
\end{center}

\newpage

%% APPENDIX E: LEHRWERK-ABGLEICH UND PRÜFUNGSBEDINGUNGEN
\begin{center}
\colorbox{spanisch}{\parbox{0.95\textwidth}{\centering\color{white}\Large\bfseries\vspace{4pt}
ANHANG E: Lehrwerk-Abgleich und Prüfungsbedingungen
\vspace{4pt}}}
\end{center}

\vspace{4mm}

\secbil{E1. Der Rahmenplan im Detail}{E1. The Curriculum Framework in Detail}

\vspace{2mm}

\begin{center}
\small
\begin{tabular}{|p{4cm}|p{10cm}|}
\hline
\rowcolor{lightgray}
\textbf{Merkmal} & \textbf{Details} \\
\hline
\textbf{Vollständiger Titel} & Rahmenplan für die Qualifikationsphase der gymnasialen Oberstufe -- Spanisch \\
\hline
\textbf{Herausgeber} & Ministerium für Bildung, Wissenschaft und Kultur des Landes Mecklenburg-Vorpommern \\
\hline
\textbf{Veröffentlichung} & August 2019 \\
\hline
\textbf{Gültig ab} & Schuljahr 2019/2020 \\
\hline
\textbf{Umfang} & 25 Seiten \\
\hline
\textbf{Jahrgangsstufen} & 10--12 (Qualifikationsphase) \\
\hline
\textbf{Gesetzliche Grundlage} & Oberstufen- und Abiturprüfungsverordnung (APVO M-V) \\
\hline
\end{tabular}
\end{center}

\vspace{3mm}

\secbil{E2. Prüfungsbedingungen für Spätbeginner}{E2. Examination Conditions for Late Starters}

\vspace{2mm}

\begin{center}
\small
\begin{tabular}{|p{5cm}|p{5cm}|p{4cm}|}
\hline
\rowcolor{lightgray}
\textbf{Regelung} & \textbf{Inhalt} & \textbf{Quelle} \\
\hline
\textbf{Wochenstunden} & 4 Wochenstunden (vierstündig) & APVO M-V \\
\hline
\textbf{Kursniveau} & Nur \textbf{grundlegendes Niveau} (kein Leistungskurs möglich) & APVO M-V \\
\hline
\textbf{Zielniveau (Gemeinsamer Europäischer Referenzrahmen)} & A2 (Ende Klasse 10) $\rightarrow$ B1(+) (Abitur) & Rahmenplan Seite 1 \\
\hline
\textbf{Prüfungsfach} & Wählbar als P4 oder P5 (mündliche Prüfung) & APVO M-V \\
\hline
\textbf{Belegpflicht} & 12 Jahreswochenstunden durchgehend bis Abitur & APVO M-V \\
\hline
\end{tabular}
\end{center}

\vspace{3mm}

\secbil{E3. Klausurdauer}{E3. Examination Duration}

\vspace{2mm}

\begin{center}
\small
\begin{tabular}{|p{5cm}|p{5cm}|p{4cm}|}
\hline
\rowcolor{lightgray}
\textbf{Phase} & \textbf{Dauer} & \textbf{Quelle} \\
\hline
\textbf{Einführungsphase} (Klasse 10) & Mindestens 45 Minuten (90 Minuten bei Aufsätzen) & APVO M-V \\
\hline
\textbf{Qualifikationsphase} (Klasse 11--12) & Mindestens 90 Minuten & APVO M-V \\
\hline
\textbf{Abitur-Semester} & Nur 1 Klausur pro Fach & APVO M-V \\
\hline
\textbf{Sprechprüfung} & Mindestens 1 Prüfung pro Schüler in der Sekundarstufe II (§ 22 APVO) & Rahmenplan Seite 19 \\
\hline
\textbf{Klausuren pro Schuljahr} & Mindestens 3 (in Fremdsprachen einschließlich neu begonnener Fremdsprache) & APVO M-V \\
\hline
\end{tabular}
\end{center}

\vspace{3mm}

\secbil{E4. Lehrwerk A\_tope.com im Überblick}{E4. Textbook A\_tope.com Overview}

\vspace{2mm}

\begin{center}
\small
\begin{tabular}{|p{4cm}|p{10cm}|}
\hline
\rowcolor{lightgray}
\textbf{Merkmal} & \textbf{Details} \\
\hline
\textbf{Titel} & A\_tope.com -- Spanisch Spätbeginner \\
\hline
\textbf{Verlag} & Cornelsen Verlag \\
\hline
\textbf{Erscheinungsjahr} & 2010 (aktualisiert 2017, Druck 2020) \\
\hline
\textbf{Zielgruppe} & Bundesweit: Gymnasium, Gesamtschule, Berufsschule, Erwachsenenbildung \\
\hline
\textbf{Zulassung} & Alle 16 Bundesländer, einschließlich Mecklenburg-Vorpommern \\
\hline
\textbf{Umfang} & 8 Unidades + 4 Module + Anhänge \\
\hline
\end{tabular}
\end{center}

\vspace{2mm}

\begin{paracol}{2}
\scriptsize
\textbf{Wichtiger Hinweis:} Das Lehrwerk A\_tope.com ist für Mecklenburg-Vorpommern \textbf{formal zugelassen}. Es wurde jedoch \textbf{vor dem Rahmenplan 2019} konzipiert (Erstausgabe 2010). Die formale Zulassung bedeutet \textbf{nicht automatisch}, dass alle verbindlichen Inhalte des M-V-Rahmenplans abgedeckt werden.

\switchcolumn

\scriptsize\color{english}
\textbf{Important note:} The textbook A\_tope.com is \textbf{formally approved} for Mecklenburg-Vorpommern. However, it was designed \textbf{before the 2019 curriculum} (first edition 2010). Formal approval does \textbf{not automatically mean} all mandatory M-V curriculum content is covered.

\end{paracol}

\vspace{3mm}

\secbil{E5. Verbindliche Inhalte laut Rahmenplan}{E5. Mandatory Content per Curriculum}

\vspace{2mm}

\begin{paracol}{2}
\scriptsize
\textbf{Rahmenplan Sek II M-V (2019), Seite 1:}

\textit{``Die in Kapitel 3.2 benannten Themen und Ziele füllen ca. \textbf{80\%} der zur Verfügung stehenden Unterrichtszeit. [...] Die linke Spalte [enthält] die \textbf{verbindlichen Inhalte}. In der rechten Spalte werden \textbf{Hinweise und Anregungen} [...] mit Beispielen zur Schwerpunktsetzung gegeben.''}

\switchcolumn

\scriptsize\color{english}
\textbf{Curriculum Sek II M-V (2019), page 1:}

\textit{``The themes [...] fill approx. \textbf{80\%} of available teaching time. [...] The left column contains \textbf{mandatory content}. The right column provides \textbf{suggestions and examples}.''}

\end{paracol}

\vspace{2mm}

\begin{center}
\scriptsize
\begin{tabular}{|c|l|c|l|}
\hline
\rowcolor{lightgray}
\textbf{Element} & \textbf{Beschreibung} & \textbf{Status} & \textbf{Quelle} \\
\hline
4 Themenfelder & El individuo, Identidad, Aspectos, Desafíos & \textbf{Verbindlich} & Rahmenplan S. 16--19 \\
\hline
Linke Spalte & z.B. ``Movimientos migratorios'' & \textbf{Verbindlich} & Rahmenplan S. 1 \\
\hline
Rechte Spalte & z.B. ``Analyse der Ursachen von Migration'' & Empfohlen & Rahmenplan S. 1 \\
\hline
Reihenfolge & Themenfeld 1 vor 2 vor 3... & Frei wählbar & Rahmenplan S. 1 \\
\hline
20\% der Zeit & Eigene Gestaltung & Frei & Rahmenplan S. 1 \\
\hline
\end{tabular}
\end{center}

\vspace{3mm}

\secbil{E6. Themenfelder 3+4: Rahmenplan versus Lehrwerk}{E6. Theme Fields 3+4: Curriculum vs. Textbook}

\vspace{2mm}

\textbf{\textcolor{spanisch}{Themenfeld 3: Aspectos actuales de la política y de la sociedad}} (30 Unterrichtsstunden Grundkurs)

\vspace{1mm}

\begin{center}
\scriptsize
\begin{tabular}{|p{4cm}|p{5cm}|p{5cm}|}
\hline
\rowcolor{lightgray}
\textbf{Verbindlicher Inhalt (Rahmenplan S. 18)} & \textbf{Hinweise im Rahmenplan} & \textbf{A\_tope.com Befund} \\
\hline
\textbf{Desarrollos económicos} & ``wirtschaftliche Entwicklungen und deren Einflüsse auf die Gesellschaft exemplarisch anhand eines oder mehrerer spanischsprachiger Länder reflektieren'' & \textcolor{spaet}{Teilweise:} Unidad 7 (Berufswelt) \\
\hline
\rowcolor{lightred}
\textbf{Movimientos migratorios} & ``Analyse und Diskussion der \textbf{Ursachen und Konsequenzen von Migration}'' [BTV] & \textcolor{spanisch}{\textbf{Nicht vorhanden}} \newline grep ``migra/inmigra/refugiad'' → 0 inhaltliche Treffer \\
\hline
\end{tabular}
\end{center}

\vspace{2mm}

\textbf{\textcolor{spanisch}{Themenfeld 4: Desafíos modernos}} (30 Unterrichtsstunden Grundkurs)

\vspace{1mm}

\begin{center}
\scriptsize
\begin{tabular}{|p{4cm}|p{5cm}|p{5cm}|}
\hline
\rowcolor{lightgray}
\textbf{Verbindlicher Inhalt (Rahmenplan S. 19)} & \textbf{Hinweise im Rahmenplan} & \textbf{A\_tope.com Befund} \\
\hline
\rowcolor{lightred}
\textbf{Tendencias ecológicas} & ``aktuelle Dokumente zum Thema \textbf{Umweltschutz} sowie Auswirkungen verschiedener Tourismuskonzepte'' [BNE] & \textcolor{spanisch}{\textbf{Nur Randerwähnung}} \newline ``turismo sostenible'' in Unidad 8 (touristische Perspektive) \newline grep ``cambio climático/contaminación'' → 0 Treffer \\
\hline
\rowcolor{lightred}
\textbf{Desafíos de la era digital} & ``Mediennutzung und Medienvielfalt im \textbf{Spannungsfeld zwischen Chancen und Risiken}'' [MD] & \textcolor{spanisch}{\textbf{Nur Vokabular}} \newline Módulo ``¿Te comunicas?'' = Handy-Aktivitäten \newline Keine kritische Reflexion \\
\hline
\end{tabular}
\end{center}

\vspace{3mm}

\secbil{E7. Direktes Lesen der Module}{E7. Direct Reading of Modules}

\vspace{2mm}

\begin{center}
\scriptsize
\begin{tabular}{|p{3.5cm}|p{2cm}|p{4cm}|p{4.5cm}|}
\hline
\rowcolor{lightgray}
\textbf{Modul} & \textbf{Seiten} & \textbf{Tatsächlicher Inhalt} & \textbf{Rahmenplan-Relevanz} \\
\hline
Vivir la diversidad & 126--127 & Toleranz-Kampagne \textbf{innerhalb Spaniens}, Decálogo gegen Diskriminierung & Diversität, \textbf{nicht} Migration \\
\hline
El comercio justo & 128--129 & 10 Gründe für fairen Handel, Kaufverhalten & ``medio ambiente'' nur als Vokabeleintrag \\
\hline
\rowcolor{lightspaet}
Una decisión importante & 130--131 & Spanischer Ingenieur \textbf{emigriert nach Deutschland} wegen Jobmangel & \textbf{Emigration} (umgekehrte Perspektive!) -- nicht Immigration nach Spanien \\
\hline
¿Te comunicas? & 132--133 & Handy-Aktivitäten: Nachrichten, Fotos, Apps & Nur Vokabular, keine kritische Medienreflexion \\
\hline
\end{tabular}
\end{center}

\vspace{3mm}

\secbil{E8. Thematische Deckung nach Themenfeld}{E8. Thematic Coverage by Theme Field}

\vspace{2mm}

\begin{center}
\small
\begin{tabular}{|c|p{4cm}|c|c|p{4cm}|}
\hline
\rowcolor{lightgray}
\textbf{Nr.} & \textbf{Themenfeld} & \textbf{Rahmenplan} & \textbf{A\_tope} & \textbf{Befund} \\
\hline
1 & El individuo y la sociedad & \checkmark & \textcolor{spaet}{\textbf{90\%}} & Unidad 1--4, 7: Jugend, Familie, Beruf \\
\hline
2 & Identidad nacional / Diversidad & \checkmark & \textcolor{spaet}{\textbf{70\%}} & Unidad 5, 6, 8: Madrid, Perú, Andalucía \newline \textit{Lücke: Nur Peru für Lateinamerika} \\
\hline
\rowcolor{lightred}
\textbf{3} & \textbf{Aspectos actuales} & \checkmark & \textcolor{spanisch}{\textbf{0\%}} & \textbf{Migration fehlt komplett} \\
\hline
\rowcolor{lightred}
\textbf{4} & \textbf{Desafíos modernos} & \checkmark & \textcolor{spanisch}{\textbf{15\%}} & Umwelt/Digital nur oberflächlich \\
\hline
\end{tabular}
\end{center}

\vspace{3mm}

\secbil{E9. Grammatik-Deckung}{E9. Grammar Coverage}

\vspace{2mm}

\begin{center}
\small
\begin{tabular}{|p{3.5cm}|c|c|p{5cm}|}
\hline
\rowcolor{lightgray}
\textbf{Grammatik} & \textbf{Niveau} & \textbf{A\_tope} & \textbf{Befund} \\
\hline
Presente, Perfecto & A2/B1 & \checkmark & Unidad 1--4 + Módulo \\
\hline
Indefinido, Imperfecto & B1 & \checkmark & Unidad 6, 8 \\
\hline
Futuro próximo & A2 & \checkmark & Unidad 4 \\
\hline
\rowcolor{lightspaet}
Subjuntivo presente & B1 & \textcolor{spaet}{Modul} & Módulo ``Vivir la diversidad'' (fakultativ) \\
\hline
\rowcolor{lightspaet}
Condicional simple & B1 & \textcolor{spaet}{Modul} & Módulo ``El comercio justo'' (fakultativ) \\
\hline
\rowcolor{lightred}
\textbf{Pluscuamperfecto} & B2 & \textcolor{spanisch}{$\times$} & grep ``pluscuam'' → 0 Treffer \\
\hline
\rowcolor{lightred}
\textbf{Voz pasiva} & B2 & \textcolor{spanisch}{$\times$} & 1 Randerwähnung bei se-Konstruktion \\
\hline
\end{tabular}
\end{center}

\vspace{2mm}

\begin{paracol}{2}
\scriptsize
\textbf{Hinweis für Spätbeginner:} Pluscuamperfecto und Passiv sind B2-Strukturen und für das Zielniveau B1(+) \textbf{nicht zwingend erforderlich}.

\switchcolumn

\scriptsize\color{english}
\textbf{Note for late starters:} Pluperfect and passive are B2 structures and \textbf{not mandatory} for the B1(+) target level.

\end{paracol}

\vspace{3mm}

\secbil{E10. Deckung nach Lernjahr}{E10. Coverage by Learning Year}

\vspace{2mm}

\begin{center}
\small
\begin{tabular}{|l|c|c|p{6cm}|}
\hline
\rowcolor{lightgray}
\textbf{Jahr / Year} & \textbf{Niveau} & \textbf{Deckung} & \textbf{Bewertung / Assessment} \\
\hline
\textbf{Jahr 1} (Klasse 10) & A2 & \textcolor{spaet}{\textbf{90\%}} & A\_tope ausreichend \\
\hline
\textbf{Jahr 2} (Klasse 11) & A2+/B1 & \textcolor{spaet}{\textbf{75\%}} & A\_tope + Module ausreichend \\
\hline
\rowcolor{lightred}
\textbf{Jahr 3} (Klasse 12) & B1(+) & \textcolor{spanisch}{\textbf{40\%}} & Ergänzung erforderlich \\
\hline
\end{tabular}
\end{center}

\vspace{3mm}

\secbil{E11. Erforderliche Ergänzungen für Jahr 3}{E11. Required Supplements for Year 3}

\vspace{2mm}

\begin{paracol}{2}
\small
\textbf{Thematisch:}
\begin{itemize}[noitemsep,topsep=2pt]
\item \textbf{Migration:} Lektüre \textit{Abdel} (Páez) + Handreichung Themenfeld 3
\item \textbf{Globalización:} Handreichung Themenfeld 3 + Aktualtexte
\item \textbf{Medio ambiente:} Handreichung Themenfeld 4 + \textit{El agua y la humanidad}
\item \textbf{Tecnología:} Aktualtexte (Módulo zu kurz)
\end{itemize}

\textbf{Grammatisch (optional für B1+):}
\begin{itemize}[noitemsep,topsep=2pt]
\item Pluscuamperfecto: externe Ergänzung
\item Passiv: externe Ergänzung
\item Subjuntivo: Módulo systematisch nutzen
\end{itemize}

\switchcolumn

\small\color{english}
\textbf{Thematically:}
\begin{itemize}[noitemsep,topsep=2pt]
\item \textbf{Migration:} Novel \textit{Abdel} (Páez) + Theme Field 3 guidelines
\item \textbf{Globalization:} Theme Field 3 guidelines + current texts
\item \textbf{Environment:} Theme Field 4 guidelines + \textit{El agua y la humanidad}
\item \textbf{Technology:} Current texts (module too short)
\end{itemize}

\textbf{Grammatically (optional for B1+):}
\begin{itemize}[noitemsep,topsep=2pt]
\item Pluperfect: external supplement
\item Passive: external supplement
\item Subjunctive: use module systematically
\end{itemize}

\end{paracol}

\vspace{3mm}

\secbil{E12. Ist Literatur verpflichtend?}{E12. Is Literature Mandatory?}

\vspace{2mm}

\begin{paracol}{2}
\small
\textbf{Antwort: NEIN}\footnotemark[18]

Das Dokument heißt ``Kommentierte Literatur\textit{empfehlungen}'' -- es ist ein Begleitdokument, keine verbindliche Vorgabe.

\textbf{ABER:} Da A\_tope das Thema \textbf{Migration nicht behandelt}, ist \textit{Abdel} die \textbf{praktisch unverzichtbare} Ergänzung für Themenfeld 3.

\textbf{\textit{Abdel} (Enrique Páez):}
\begin{itemize}[noitemsep,topsep=2pt]
\item 16-jähriger Protagonist aus Marokko
\item Thema: Illegale Immigration
\item Niveau: A2/B1 (sprachlich zugänglich)
\item Ideal für Spätbeginner
\end{itemize}

\switchcolumn

\small\color{english}
\textbf{Answer: NO}\footnotemark[18]

The document is titled ``Annotated Literature \textit{Recommendations}'' -- it's a supplementary document, not a binding requirement.

\textbf{BUT:} Since A\_tope \textbf{does not cover migration}, \textit{Abdel} is the \textbf{practically indispensable} supplement for theme field 3.

\textbf{\textit{Abdel} (Enrique Páez):}
\begin{itemize}[noitemsep,topsep=2pt]
\item 16-year-old protagonist from Morocco
\item Topic: Illegal immigration
\item Level: A2/B1 (linguistically accessible)
\item Ideal for late starters
\end{itemize}

\end{paracol}

\footnotetext[18]{Literaturempfehlungen Spanisch M-V, Zeilen 113--118: \textit{Abdel} als Empfehlung, nicht Pflicht. RP Sek II (2019), S. 1: ``80\% der Unterrichtszeit [...] Freiraum für eigene Unterrichtsgestaltung.''}

\vspace{3mm}

\newpage

\secbil{E13. Kurszusammensetzung}{E13. Course Composition}

\vspace{2mm}

\begin{center}
\small
\begin{tabular}{|c|c|p{10cm}|}
\hline
\rowcolor{lightgray}
\textbf{Klasse} & \textbf{Schüler} & \textbf{Besonderheiten} \\
\hline
\textbf{10} & 6 & Regelkurs, keine besonderen Förderbedarfe \\
\hline
\textbf{11} & 2 & \textbf{Schülerin A:} Arbeitet mit Klasse 10 (Auslandsaufenthalt, muss Stoff nachholen) \newline \textbf{Schülerin B:} Sehr begabt, arbeitet praktisch mit Klasse 12 \\
\hline
\textbf{12} & 2 & \textbf{Schülerin A:} Lese-Rechtschreib-Schwäche mit Nachteilsausgleich, schwächeres Sprachniveau, bemüht \newline \textbf{Schüler B:} Mittelmäßiges Niveau, inkonsistente Mitarbeit \\
\hline
\rowcolor{lightspaet}
\textbf{Gesamt} & \textbf{10} & Jahrgangsübergreifend, starke Binnendifferenzierung erforderlich \\
\hline
\end{tabular}
\end{center}

\vspace{2mm}

\begin{paracol}{2}
\scriptsize
\textbf{Konsequenzen für die Unterrichtsplanung:}
\begin{itemize}[noitemsep,topsep=2pt]
\item Effektiv 3 Lerngruppen in einem Kurs
\item Klasse 10: Grundlagenvermittlung (A1$\rightarrow$A2)
\item Klasse 11/12: Vertiefung (A2$\rightarrow$B1)
\item Differenzierte Materialien erforderlich
\item Nachteilsausgleich für Lese-Rechtschreib-Schwäche beachten
\end{itemize}

\switchcolumn

\scriptsize\color{english}
\textbf{Implications for lesson planning:}
\begin{itemize}[noitemsep,topsep=2pt]
\item Effectively 3 learning groups in one course
\item Grade 10: Foundation (A1$\rightarrow$A2)
\item Grade 11/12: Deepening (A2$\rightarrow$B1)
\item Differentiated materials required
\item Accommodate reading/writing difficulties
\end{itemize}

\end{paracol}

\vspace{3mm}

\secbil{E14. Gesamtfazit Lehrwerk-Eignung}{E14. Overall Assessment of Textbook Suitability}

\vspace{2mm}

\begin{center}
\fcolorbox{spanisch}{lightred}{\parbox{0.9\textwidth}{
\centering\small
\textbf{Kernproblem: A\_tope.com (2010/2017) wurde VOR dem Rahmenplan M-V (2019) konzipiert.}\\[6pt]
Die \textbf{formale Zulassung} für M-V (alle 16 Bundesländer) bedeutet \textbf{nicht}, dass alle\\
\textbf{verbindlichen Inhalte} des M-V-Rahmenplans abgedeckt werden.
}}
\end{center}

\vspace{3mm}

\begin{paracol}{2}
\small
\textbf{Frage 1: Kann die Spätbeginner-Sequenz komplett auf A\_tope aufbauen?}

\textcolor{spanisch}{\textbf{NEIN}} -- Für die verbindlichen Themenfelder 3+4 fehlen:
\begin{itemize}[noitemsep,topsep=2pt]
\item \textbf{Movimientos migratorios} (0 Treffer)
\item \textbf{Tendencias ecológicas} (nur ``turismo sostenible'')
\item \textbf{Desafíos de la era digital} (nur Vokabular)
\end{itemize}

\textbf{Frage 2: Sind Handreichungen erforderlich?}

\textcolor{spanisch}{\textbf{JA}} -- Für Klasse 12 zwingend:
\begin{itemize}[noitemsep,topsep=2pt]
\item Handreichung Themenfeld 3 (Migration, Globalización)
\item Handreichung Themenfeld 4 (Umwelt, Digitalisierung)
\item Lektüre \textit{Abdel} (Enrique Páez)
\end{itemize}

\switchcolumn

\small\color{english}
\textbf{Question 1: Can the late starter sequence be based entirely on A\_tope?}

\textcolor{spanisch}{\textbf{NO}} -- Missing for mandatory theme fields 3+4:
\begin{itemize}[noitemsep,topsep=2pt]
\item \textbf{Migration movements} (0 hits)
\item \textbf{Ecological trends} (only ``sustainable tourism'')
\item \textbf{Digital challenges} (vocabulary only)
\end{itemize}

\textbf{Question 2: Are supplementary guidelines required?}

\textcolor{spanisch}{\textbf{YES}} -- Mandatory for grade 12:
\begin{itemize}[noitemsep,topsep=2pt]
\item Guidelines Theme Field 3 (Migration, Globalization)
\item Guidelines Theme Field 4 (Environment, Digitalization)
\item Novel \textit{Abdel} (Enrique Páez)
\end{itemize}

\end{paracol}

\vspace{3mm}

\begin{center}
\fcolorbox{spaet}{lightspaet}{\parbox{0.9\textwidth}{
\centering\small
\textbf{Empfehlung / Recommendation:}\\[4pt]
A\_tope.com als \textbf{Basis für Jahr 1--2} (Klasse 10--11) nutzen.\\
Für \textbf{Jahr 3} (Klasse 12) \textbf{zwingend ergänzen} mit Handreichungen + Lektüre.\\[4pt]
\scriptsize\textit{Prüfung: Januar 2026 | Methode: Direktes Lesen + Grep gegen OCR-Volltext (1,5 MB, 10.240 Zeilen)}
}}
\end{center}

\vspace{6mm}

%% ABBREVIATIONS FINAL
\vspace{4mm}
\secbil{Abkürzungsverzeichnis}{Abbreviations}

\vspace{2mm}

\begin{center}
\scriptsize
\begin{tabular}{|l|p{5cm}|p{5cm}|}
\hline
\rowcolor{lightgray}
\textbf{Abk.} & \textbf{Deutsch} & \textbf{English} \\
\hline
GER/GeR & Gemeinsamer Europäischer Referenzrahmen & Common European Framework of Reference (CEFR) \\
\hline
\textbf{APVO M-V} & \textbf{Oberstufen- und Abiturprüfungsverordnung} & \textbf{Upper Secondary \& Abitur Exam Regulations} \\
\hline
GK & Grundkurs (fortgef. FS: 3 WStd) & Basic Course (continued FL: 3 hrs/wk) \\
\hline
\textcolor{spaet}{\textbf{Neu beg. FS}} & \textcolor{spaet}{\textbf{Spätbeginner: 4 WStd}} & \textcolor{spaet}{\textbf{Late starter: 4 hrs/wk}} \\
\hline
LK & Leistungskurs (5 WStd) & Advanced Course (5 hrs/wk) \\
\hline
AFB & Anforderungsbereich (I/II/III) & Requirement Level (I/II/III) \\
\hline
RP & Rahmenplan (2019) & Curriculum Framework (2019) \\
\hline
HR-SP & Handreichung Sprechprüfung (2020) & Oral Exam Guidelines (2020) \\
\hline
HR-PP & Handreichung Paarprüfung (2025) & Pair Exam Guidelines (2025) \\
\hline
VV & Verwaltungsvorschrift & Administrative Regulation \\
\hline
P4/P5 & Prüfungsfach 4/5 (Abitur) & Exam subject 4/5 (Abitur) \\
\hline
SuS & Schülerinnen und Schüler & Students \\
\hline
TF & Themenfeld & Theme Field \\
\hline
HR & Handreichung & Guidelines \\
\hline
\end{tabular}
\end{center}

\vspace{4mm}

%% SOURCES
\begin{center}
\fcolorbox{spaet}{lightspaet}{\parbox{0.95\textwidth}{
\centering
\small
\textbf{Quellenverzeichnis / Sources:}\\[2pt]
\scriptsize
$^{1,4,6}$ RP Spanisch Qualifikationsphase M-V (2019) |
$^{3}$ \textbf{APVO M-V (2019): Stundentafel} |
$^{5,9}$ HR-SP: Handreichung Sprechprüfung (2020) |\\
$^{7}$ HR-PP: Handreichung Paarprüfung Abitur (2025) |
$^{8}$ Kommentierte Literaturempfehlungen |
$^{10}$ Grammatik = Synthese |\\
$^{11}$ VV Sprechprüfung (24.07.2020) |
$^{12}$ HR-SP Anlage II: Bewertungsbögen |
$^{13}$ VV Paarprüfung (03.12.2025) |
$^{14}$ Vergleich HR-SP/HR-PP |\\
$^{15}$ APVO: Spätbeginner = 4 WStd |
$^{16}$ HR-SP Timing-Empfehlung |
$^{17}$ Jahrgangsübergreifend |
$^{18}$ Literaturempfehlungen (Begleitdokument)\\[4pt]
\textit{Hinweis: Wochensequenzen und Grammatik-Progression sind didaktische Synthesen.}\\[4pt]
\textcolor{darkgray}{Lernpfade | Spanisch Spätbeginner | M-V | v1.9 | Januar 2026}
}}
\end{center}

\end{document}
